\section{Literatür verisiyle yazılım uygulamaları}
\label{sec:spss-b06}

\textbf{Amaç ve veri.} \texttt{b06.csv} içinde 395 kayıt, yalnız bu
uygulama için sonlu bir örnekleme çerçevesidir \citep{cortez2008data}.
\texttt{srs40} geri koymadan basit rastgele seçilen 40 kaydı;
\texttt{strat40} GP'den 35, MS'den 5 kayıt seçimini gösterir. Seçimler sabit
tohumla önceden üretilmiştir.

\textbf{Ortak veri.} Doğrulanan B06 uygulamasında üç yazılım da
\nolinkurl{companion/spss/csv/b06.csv} dosyasını kullanmıştır. Dosya
395 satır ve 13 sütun içerir; B01'in öğrenci değişkenlerine seçim
göstergeleri ve \texttt{strw} ağırlığı eklenmiştir. B01 dosyasını yeniden
adlandırmak yeterli değildir. Seçimler yeniden üretilmez; aynı hazır
göstergeler kullanılır. Bu karşılaştırma Excel içe aktarmanın doğrulandığı
anlamına gelmez. Kodları \texttt{companion} klasörünü içeren paket kökünden
çalıştırın; veri sözlüğü ve hazırlık için bkz. sayfa~\pageref{sec:spss-guide}.

\subsection{IBM SPSS uygulaması}
\textbf{Başlamadan önce.} Aşağıdaki komut dosyası B06 CSV verisini
açar ve analiz eder. Menü yolunu kullanacaksanız önce veri açma ve
tanımlama komutlarını \texttt{EXECUTE} satırı dahil çalıştırın.
Dosya açma hatası alırsanız \texttt{/FILE} yolunu ve dosyanın gerçekten
bu konumda bulunduğunu kontrol edin; veri açılmadan sonraki komutlar
çalışmaz.

\begin{spssadimlar}
\spssadim{Henüz filtre ve ağırlık yokken \emph{Analyze $\to$ Descriptive Statistics $\to$ Descriptives} yolunda \texttt{g3} için ortalama ve standart sapmayı alın; tam dosyada 395 kayıt bulunduğunu kontrol edin.}
\spssadim{\emph{Data $\to$ Select Cases} yolunda \emph{If condition is satisfied $\to$ If} seçin; \texttt{srs40=1} yazın. \emph{Continue} ardından \emph{Filter out unselected cases} ve \emph{OK} seçin; kayıtları silmeyin.}
\spssadim{\emph{Analyze $\to$ Descriptive Statistics $\to$ Descriptives} yolunda \texttt{g3} için ortalama ve standart sapmayı alın; analizde 40 kayıt kullanıldığını kontrol edin.}
\spssadim{\emph{Data $\to$ Select Cases} penceresinde önce \emph{All cases} ile filtreyi kapatın. Ardından koşulu \texttt{strat40=1} yaparak yeniden filtreleyin.}
\spssadim{\emph{Analyze $\to$ Descriptive Statistics $\to$ Frequencies} yolunda \texttt{school} tablosunu alın. Henüz ağırlık vermeden GP=35, MS=5 sayımlarını kontrol edin.}
\spssadim{\emph{Data $\to$ Weight Cases} yolunda \emph{Weight cases by} seçin, \texttt{strw} değişkenini ağırlık alanına taşıyıp \emph{OK} seçin. \emph{Descriptives} ile \texttt{g3} ortalamasını tekrar alın.}
\spssadim{Bitirirken \emph{Data $\to$ Weight Cases $\to$ Do not weight cases} ve \emph{Data $\to$ Select Cases $\to$ All cases} seçin. Hem menü hem aşağıdaki komut yolu 395 kaydı korur; kayıtları kalıcı olarak silmeyin.}
\end{spssadimlar}

{\raggedright
\noindent\textbf{Komutların görevi.} \texttt{FILTER BY} seçili kayıtları analiz eder, diğerlerini silmez. \texttt{FILTER OFF} tüm kayıtları yeniden etkinleştirir. \texttt{WEIGHT BY} ağırlıklı nokta tahminini hesaplatır; işlem sonunda \texttt{WEIGHT OFF} ile kapatılır.\par}

\textbf{Uygulama kodu:} \texttt{b06}. Doğrulanan veri dosyası
\texttt{companion/spss} altındaki \texttt{csv/b06.csv} dosyasıdır.

\begin{lstlisting}[style=prismSPSS]
SET DECIMAL=DOT.
GET DATA /TYPE=TXT
 /FILE='companion/spss/csv/b06.csv'
 /ENCODING='UTF8'
 /ARRANGEMENT=DELIMITED /DELCASE=LINE
 /DELIMITERS="," /QUALIFIER='"'
 /FIRSTCASE=2
 /VARIABLES=id F8.0 school F8.0 sex F8.0 age F8.0
 studytime F8.0 absences F8.0 g1 F8.0 g2 F8.0
 g3 F8.0 pass10 F8.0 srs40 F8.0 strat40 F8.0
 strw F20.12.
FILTER OFF.
WEIGHT OFF.
SPLIT FILE OFF.
VARIABLE LABELS
 g3 'Yil sonu matematik notu'
 srs40 'Basit rastgele orneklem secimi'
 strat40 'Tabakali orneklem secimi'
 strw 'Tabaka agirligi'.
VALUE LABELS school 1 'GP' 2 'MS'.
VALUE LABELS srs40 strat40 0 'Secilmedi' 1 'Secildi'.
VARIABLE LEVEL
 id school sex pass10 srs40 strat40 (NOMINAL)
 studytime (ORDINAL)
 age absences g1 g2 g3 strw (SCALE).
FORMATS g3 strw (F12.6).
EXECUTE.
DESCRIPTIVES VARIABLES=g3 /STATISTICS=MEAN STDDEV.
FILTER BY srs40.
DESCRIPTIVES VARIABLES=g3 /STATISTICS=MEAN STDDEV.
FILTER OFF.
FILTER BY strat40.
FREQUENCIES VARIABLES=school.
WEIGHT BY strw.
DESCRIPTIVES VARIABLES=g3 /STATISTICS=MEAN.
WEIGHT OFF.
FILTER OFF.
EXECUTE.
\end{lstlisting}

\begin{outputguidebox}
\textbf{Paylaşılan SPSS çıktısıyla doğrulanan değerler.}
Basit rastgele seçilen 40 kaydın not ortalaması \SPSrsMean'dir.
Tabakalı seçimde GP için $349/35$, MS için $46/5$ ağırlıkları kullanıldığında
ortalama \SPStratMean\ bulunur. Karşılaştırma değeri, dosyanın tamamının
\SPMatMean\ ortalamasıdır. Tek bir çekimde hangisinin daha yakın olduğu,
yöntemlerden birinin her zaman daha iyi olduğunu kanıtlamaz.
Ağırlıklı çıktıda gösterilen toplam ağırlık 395'tir; gerçek seçilmiş
öğrenci sayısı yine 40'tır. \texttt{WEIGHT BY} burada yalnız ağırlıklı nokta
tahminini göstermek içindir, karmaşık örnekleme standart hatası üretmez.
\end{outputguidebox}


\par\noindent\begin{minipage}{\linewidth}
\subsection{R uygulaması}
\label{sec:r-gercek-b06}

\texttt{subset} seçili kayıtları ayrı nesneye alır; \texttt{weighted.mean} tabaka ağırlıklı ortalamayı verir.

\begin{lstlisting}[style=prismR]
veri <- read.csv("companion/spss/csv/b06.csv")
stopifnot(nrow(veri) == 395, ncol(veri) == 13,
          !anyNA(veri))
basit <- subset(veri, srs40 == 1)
tabakali <- subset(veri, strat40 == 1)
stopifnot(nrow(basit) == 40, nrow(tabakali) == 40)
ozet <- function(degerler)
  c(n=length(degerler), mean=mean(degerler),
    sd=sd(degerler))
print(cbind(tam_dosya=ozet(veri$g3),
            basit_rastgele=ozet(basit$g3)))
okul <- factor(tabakali$school, c(1, 2), c("GP", "MS"))
frekans <- table(okul)
print(frekans); print(100 * prop.table(frekans))
print(nrow(tabakali)); print(sum(tabakali$strw))
print(weighted.mean(tabakali$g3, tabakali$strw))
\end{lstlisting}

\textbf{Çıktıyı okuma.} Paylaşılan R çıktısı 395 satır, 13 sütun ve
her sütunda sıfır eksik değer gösterir. \texttt{n}, \texttt{mean} ve
\texttt{sd} satırları tam dosya ile basit rastgele örneklemi karşılaştırır.
Tabakalı örneklemde okul frekansları ağırlıksızdır; 40 kayıt ile 395
ağırlık toplamı ayrı yazdırılır. Ana \texttt{veri} nesnesi değişmez.
\end{minipage}\par

\par\noindent\begin{minipage}{\linewidth}
\subsection{Python uygulaması}
\label{sec:python-b06}

\texttt{loc} ve \texttt{copy} seçimi ayrı veri çerçevesinde tutar. \texttt{np.average} ağırlıklı ortalamayı hesaplar.

\begin{lstlisting}[style=prismPythonApp]
import pandas as pd
import numpy as np

veri = pd.read_csv("companion/spss/csv/b06.csv")
assert veri.shape == (395, 13)
assert not veri.isna().any().any()
basit = veri.loc[veri.srs40 == 1].copy()
tabakali = veri.loc[veri.strat40 == 1].copy()
assert len(basit) == len(tabakali) == 40
print(pd.DataFrame({
    "tam_dosya": veri.g3.agg(["count", "mean", "std"]),
    "basit_rastgele": basit.g3.agg(["count", "mean", "std"])
}))
okul = tabakali.school.map({1: "GP", 2: "MS"})
frekans = okul.value_counts().reindex(["GP", "MS"])
print(frekans); print(100 * frekans / len(tabakali))
print(len(tabakali)); print(tabakali.strw.sum())
print(np.average(tabakali.g3, weights=tabakali.strw))
\end{lstlisting}

\textbf{Çıktıyı okuma.} Paylaşılan Python çıktısı da 395 satır,
13 sütun ve her sütunda sıfır eksik değer gösterir. \texttt{count}, R'deki
\texttt{n}; \texttt{std} ise \texttt{sd} satırına karşılık gelir.
Okul frekansları, yüzdeler ve ağırlıklı ortalama R ve SPSS ile uyumludur.
Bu uygulama bir hipotez testi veya $p$ değeri üretmez; ana veri korunur.
\end{minipage}\par

\subsection{Sonuçların karşılaştırılması ve ortak rapor}
12 Eylül 2026'da paylaşılan SPSS tabloları ile R ve Python metin
çıktıları karşılaştırılmıştır. Tam dosya ve basit rastgele örneklemin
gözlem sayısı, ortalaması ve standart sapmasından oluşan altı hücre ile
tabakalı seçimin ağırlıklı ortalaması, gösterilen yuvarlama
hassasiyetlerinde uyumludur. Okul frekansları, yüzdeler, tabakalı
örneklem büyüklüğü ve ağırlık toplamı da eşleşmiştir. Proje CSV'sinden
yapılan hesaplar bu değerleri destekler. Aşağıdaki tablolar çıktıların
kitap düzenine uyarlanmış özetidir; ekran görüntüsü değildir.

\begin{table}[H]
\centering
\caption{B06'da tam dosya ve seçilmiş örneklemlerin not özetleri}
\label{tab:b06-dogrulanmis-betimsel}
\small
\begin{tabular}{@{}lrrr@{}}
\toprule
İncelenen veri & Kayıt sayısı & Ortalama & Std. sapma \\
\midrule
Tam dosya & 395 & 10,415190 & 4,581443 \\
Basit rastgele örneklem & 40 & 10,275000 & 5,565404 \\
Tabakalı örneklem (ağırlıklı) & 40 & 10,456347 & --- \\
\bottomrule
\end{tabular}
\par\smallskip
{\footnotesize Ondalıklı değerler altı basamağa yuvarlanmıştır.
İlk iki satırın standart sapmaları sırasıyla 394 ve 39 böleniyle
hesaplanmıştır; standart hata değildir. Son satırda yalnız ağırlıklı
ortalama raporlanmıştır; çizgi sıfır anlamına gelmez. SPSS'in bu
satır için gösterdiği $N=395$ ağırlık toplamıdır, gerçek kayıt sayısı
40'tır.\par}
\end{table}

\begin{table}[H]
\centering
\caption{B06 tabakalı örnekleminde ağırlıksız okul frekansları}
\label{tab:b06-dogrulanmis-tabaka}
\begin{tabular}{@{}lrr@{}}
\toprule
Okul & Frekans & Yüzde \\
\midrule
GP & 35 & 87,5 \\
MS & 5 & 12,5 \\
Toplam & 40 & 100,0 \\
\bottomrule
\end{tabular}
\par\smallskip
{\footnotesize Yüzdelerin paydası seçilen 40 kayıttır. Okul değişkeninde
eksik değer olmadığından SPSS'in \emph{Percent} ve \emph{Valid Percent}
sütunları eşittir. Ağırlıklar bu tablo alındıktan sonra etkinleştirilir.\par}
\end{table}

\Needspace{15\baselineskip}
\begin{outputguidebox}
Tablo~\ref{tab:b06-dogrulanmis-betimsel} içindeki tam dosya ortalaması,
yalnız bu 395 kayıt eğitim amaçlı sonlu evren kabul edildiğinde bilinen
evren ortalamasıdır. Basit rastgele örneklemin 10,275000 ve tabakalı
örneklemin 10,456347 ortalamaları bu değerin iki ayrı tahminidir.
Bu tek seçimde tabakalı tahmin daha yakındır; bu durum yöntemin her
örneklemde daha yakın sonuç vereceğini kanıtlamaz.

Tablo~\ref{tab:b06-dogrulanmis-tabaka} içinde GP için 35, MS için 5
kayıt vardır. Ağırlıklar sırasıyla $349/35$ ve $46/5$ olduğundan
toplamları $35(349/35)+5(46/5)=395$ olur. Bu işlem 355 yeni öğrenci
eklemez. Ağırlıklı nokta tahmini, tasarıma uygun standart hata veya
güven aralığı hesabı değildir.

Seçim göstergeleri öğretim amacıyla hazırlanmıştır; kaynak araştırmanın
özgün örnekleme tasarımını göstermez. Yazılımların hesaplama uyumu,
395 kaydın daha geniş bir öğrenci evrenini temsil ettiğini kanıtlamaz.
\end{outputguidebox}

\textbf{Raporlama örneği (doğrulanmış çıktılarla).}
``395 kayıtlık öğretim çerçevesinin not ortalaması 10,415190'dır.
Basit rastgele seçilmiş 40 kaydın ortalaması 10,275000, standart sapması
5,565404 bulunmuştur. GP'den 35 ve MS'den 5 kayıt içeren ayrı tabakalı
örneklemde, tabaka ağırlıklarıyla hesaplanan ortalama 10,456347'dir.
Ağırlıkların toplamı 395 olmakla birlikte gerçek örneklem büyüklüğü
40'tır. Karşılaştırma tek seçime aittir; yöntemlerin genel üstünlüğü
veya daha geniş bir evrene genellenebilirlik iddiası taşımamaktadır.''


\textbf{Görev.} $\bar x_{\rm tab}=(349/395)\bar x_{GP}+(46/395)\bar x_{MS}$
ile sonucu elle doğrulayın. Bu seçimlerin kaynak araştırmanın özgün
örnekleme tasarımı olmadığını raporlayın. Sonraki uygulamaya geçmeden
filtre ve ağırlığın kapalı olduğunu, etkin dosyanın 395 kaydı koruduğunu
kontrol edin.